\chapter{其他约化群的Langlands对应II：Satake同构}\label{L19}\lecturedate{2020/04/22}\sourcepage{157}
\section{约化群的Bruhat分解}\label{L19S1}\idx{Bruhat分解}\idx{抛物子群}
设$k$为域，$G$为连通约化$k$群。$G$的\textbf{抛物子群}$P$是使$G/P$为射影簇的闭子群；等价地，在代数闭包上，$P$包含$G$的一个Borel子群。一般而言，$G$在$k$上未必拟分裂。因此，我们使用$G$的极小抛物$k$子群来代替Borel子群。任意两个极小抛物$k$子群都在$k$上共轭。事实上，若两个抛物$k$子群在$k$的某个扩域上共轭，则它们已经在$k$上共轭。

设$P$为$G$的极小抛物$k$子群。存在$G$的极大$k$分裂环面$S$，使$P$为$k$群的半直积$P=Z(S)\ltimes R_u(P)$，其中$Z(S)$表示$S$的中心化子，$R_u(P)$为$P$的幂幺根基，因而是正规子群。如\lref{L18}所述，${}_k\Phi=\Phi(G,S)$为相对根系，${}_kW$为$G$的相对Weyl群。此外，${}_kW\cong N(S)/Z(S)$，其中$N(S)$为$S$在$G$中的正规化子。$G$的Lie代数$\g$也有分解
\begin{equation}\tag{1.1}\label{L19eq11}\g=\g_0\oplus\bigoplus_{\alpha\in{}_k\Phi}\g_\alpha.\end{equation}
每个相对根$\alpha\in{}_k\Phi$都有相应的、由$S$正规化的幂幺$k$根子群$U_\alpha$；其Lie代数为$\bigoplus_{n\ge1}\g_{n\alpha}$，其中只取仍为根的$n\alpha$。可以给${}_k\Phi$选定一个正负划分，使$P$由$Z(S)$及正根对应的$U_\alpha$生成。记这一划分确定的基为${}_k\Delta$。对任意子集$\theta\subseteq{}_k\Delta$，以$P_\theta$表示$G$中由$Z(S)$及下述$U_\alpha$生成的子群：把$\alpha\in{}_k\Phi$写成单根的线性组合后，不属于$\theta$的单根的系数均非负。
% [Tutor Note: 原书 p157 作 $U_\alpha$ 的 Lie 代数为 $\g_\alpha$；相对根系若非约化（如 $BC_n$ 型），根子群 $U_\alpha$ 的 Lie 代数应为 $\bigoplus_{n\ge1}\g_{n\alpha}$（只取仍为根的 $n\alpha$），译文按后者处理。]
% [Tutor Note: 原书 p157–158 称 $P_\theta$ 是 $P$ 的子群；但 $\theta={}_k\Delta$ 时 $P_{{}_k\Delta}=G$ 显然不在 $P$ 中，故 $P_\theta$ 应理解为 $G$ 中由 $Z(S)$ 与这些 $U_\alpha$ 生成的子群，译文已改。]
\sourcepage{158}
于是$P_{{}_k\Delta}=G$、$P_\varnothing=P$；若$\theta_1\subseteq\theta_2$，则$P_{\theta_1}\subseteq P_{\theta_2}$。这些$P_\theta$称为包含$P$的$G$的\textbf{标准抛物子群}。$G$的任意抛物$k$子群都在$k$上共轭于唯一一个$P_\theta$。

令$U:=R_u(P)$。由${}_kW\cong N(S)/Z(S)$，对每个$w\in{}_kW$，可在$N(S)(k)$中选取代表元$n_w$。$G(k)$的\textbf{Bruhat分解}\znote{Bruhat 分解把 $G(k)$ 按双陪集 $U(k)n_wP(k)$ 拆成 $|{}_kW|$ 块，每块对应一个相对 Weyl 群元素。取 $k=\bar k$、$P=B$ 时，$Un_wB/B$ 是旗簇 $\mathcal B$ 中的仿射胞腔 $\A^{\ell(w)}$，其闭包给出 Bruhat 序，这是 Schubert 演算与相交上同调的起点。对 $p$ 进群，同一分解又是球面 Hecke 代数基 $c_\lambda$ 与 Satake 同构的组合基础。注意 $P$ 不必是 Borel、$U$ 也不必是极大幂幺子群，故该分解对任意抛物子群都成立，这使它能直接用于约化群的相对理论。[来源：nLab "Bruhat decomposition"，\url{https://ncatlab.org/nlab/show/Bruhat+decomposition}（Last revised 2024-07-11；2026-11-14 访问）。]}断言，$G(k)$为双陪集$U(k)n_wP(k)$的不交并，即
\begin{equation}\tag{1.2}\label{L19eq12}G(k)=\bigsqcup_{w\in{}_kW}U(k)n_wP(k).\end{equation}
按定义，$w(\alpha)(s)=\alpha(n_w^{-1}sn_w)$，$s\in S$，从而$n_w^{-1}U_\alpha n_w=U_{w^{-1}(\alpha)}$。
% [Tutor Note: 原书 p158 作 $n_w^{-1}U_\alpha n_w=U_{w(\alpha)}$；按本句给出的 $w(\alpha)(s)=\alpha(n_w^{-1}sn_w)$，正确关系为 $n_w U_\alpha n_w^{-1}=U_{w(\alpha)}$，等价地 $n_w^{-1}U_\alpha n_w=U_{w^{-1}(\alpha)}$，译文按正确写法处理。]定义
\[
U'_w:=\prod_{\substack{\alpha>0,\ w^{-1}(\alpha)<0\\\alpha/2\notin{}_k\Phi}}U_\alpha,
\qquad
U''_w:=\prod_{\substack{\alpha>0,\ w^{-1}(\alpha)>0\\\alpha/2\notin{}_k\Phi}}U_\alpha.
\]
乘积按适当顺序取。乘法给出簇的同构$U'_w\times U''_w\cong U$，并且$U''_w(k)n_wP(k)=n_wP(k)$。映射$U'_w(k)\times P(k)\to U(k)n_wP(k)$，$(u,p)\mapsto un_wp$，为双射。因此
% [Tutor Note: 原书 p158 定义 $U'_w,U''_w$ 时用条件 $w(\alpha)<0$、$w(\alpha)>0$ 且无 $\alpha/2$ 条件；译文改用 $w^{-1}(\alpha)$ 比较（与上式修正后的共轭关系一致），并补上 $\alpha/2\notin{}_k\Phi$（相对根系非约化时避免重复计算 $\g_{n\alpha}$）。]
\begin{equation}\tag{1.3}\label{L19eq13}(G/P)(k)=G(k)/P(k)=\bigsqcup_{w\in{}_kW}U'_w(k)n_wP(k)/P(k).\end{equation}
% [Tutor Note: 原书 p158 的 (1.3) 右端只写到 $U'_w(k)n_w$，漏了 $P(k)$ 与商 $/P(k)$；按上一段的双射，$U'_w(k)n_wP(k)/P(k)$ 才是 $G(k)/P(k)$ 的分解，译文已补。]
还具有以下性质。
\begin{itemize}
\item 若$k=\bar k$，则$P=B$为Borel子群，商$\mathcal B:=G/B$为$G$的\textbf{旗簇}\znote{旗簇 $\mathcal B=G/B$ 参数化 $G$ 的所有 Borel 子群（即「完全旗」）；$G=\GL_n$ 时它就是通常的完全旗流形，是射影簇。取定 $B$ 后，$\mathcal B$ 被 Bruhat 胞腔 $\mathcal B_w=Un_wB/B\cong\A^{\ell(w)}$ 分解，闭包 $\overline{\mathcal B_w}$ 是 Schubert 簇，维数等于 $w$ 的长度。它也是几何表示论中「$G$ 作用在自身 Borel 子群上」的紧化，几何 Satake 等价把它的等变层范畴与 $G^\vee$ 的表示联系起来。[来源：nLab "flag variety"，\url{https://ncatlab.org/nlab/show/flag+variety}（Last revised 2026-09-28；2026-11-14 访问）。]}。令$\mathcal B_w=U'_wn_wB/B=Un_wB/B$，则
\[\mathcal B_w\cong\A_k^{\ell(w)},\]
其中$\ell(w)$为$w$的长度：$w$可写成$\ell(w)$个单根反射的乘积，而任何这样的表达式都至少需要$\ell(w)$个反射。因此，此时\eqref{L19eq13}给出$\mathcal B$的胞腔分解。设$\overline{\mathcal B_w}$为$\mathcal B_w$的闭包，则
\[\overline{\mathcal B_w}=\bigsqcup_{y\preceq w}\mathcal B_y,\]
其中$\preceq$为\textbf{Bruhat序}，由关系
\[y\prec s_\alpha y\quad\text{当}\quad\ell(s_\alpha y)=\ell(y)+1\]
生成，这里$\alpha$可取任意根。于是，$y\preceq w$当且仅当$\ell(y)\le\ell(w)$，并且$w$的某个约化表达式$w=s_{\alpha_{\ell(w)}}\cdots s_{\alpha_1}$中存在一个子词，其乘积为$y$，该子词也可以取成$y$的约化表达式。
% [Tutor Note: 原书 p158 末作「存在 $w$ 的约化表达式使 $y=s_{\alpha_{\ell(y)}}\cdots s_{\alpha_1}$」，即要求 $y$ 是 $w$ 的约化表达式的**前缀**；该判据一般不对（如 $y$ 可以是中段子词），正确说法是「$y$ 的某个约化表达式是 $w$ 的约化表达式的子词」，译文按子词表述。]
\sourcepage{159}
\item 设$\theta\subseteq{}_k\Delta$，以$W_\theta$表示由$s_\alpha$，$\alpha\in\theta$，生成的${}_kW$的子群。对${}_k\Delta$的任意两个子集$\theta,\theta'$，有双陪集之间的双射
\[P_\theta(k)\backslash G(k)/P_{\theta'}(k)\cong W_\theta\backslash{}_kW/W_{\theta'}.\]
\end{itemize}\idx{旗簇}\idx{Bruhat序}
\originalnumber{1}
\begin{example}\label{L19S1O1}
我们给出一个非分裂环面的例子。定义
\[T_{\R}:=\{(x,y)\in\R^2\mid x^2+y^2=1\},\]
这是在$\R$上不分裂的一维环面。为了说明它是环面，可构造同构
\[f:T_{\R}\otimes_{\R}\C\xrightarrow{\sim}\C^\times,\qquad(x,y)\longmapsto x+yi.\]
注意，$T_{\R}\otimes_{\R}\C=\{(x,y)\in\C^2\mid x^2+y^2=1\}$，而逆映射$f^{-1}$为
\[\C^\times\ni z\longmapsto\left(\frac{z^2+1}{2z},\frac{z^2-1}{2iz}\right)\in T_{\R}\otimes_{\R}\C.\]
因此$T_{\R}\otimes_{\R}\C\cong\C^\times$。然而，它在$\R$上显然不同构于分裂环面$\mathbb G_m$：$T_{\R}(\R)$有界且紧，而$\R^\times$并不紧。
\end{example}
\section{球面Hecke代数}\label{L19S2}\idx{球面Hecke代数}\idx{卷积}
现在设$F$为非阿基米德局部域，$\OO_F$为$F$的整数环，$\mathfrak m\subset\OO_F$为极大理想，$\pi\in\mathfrak m$为素元。设$G$为$\OO_F$上的分裂约化群，即它具有一个在$\OO_F$上分裂的极大环面$T$，使$T(\OO_F)\cong(\OO_F^\times)^r$，其中$r$为$G$的秩。在$G(F)$上存在Haar测度$\mu$，可归一化为$\mu(G(\OO_F))=1$。

以$C_c(G)$表示\textbf{Schwartz函数}空间，即$G(F)$上具有紧支撑的局部常值复值函数所组成的空间。对任意两个Schwartz函数$f_1,f_2$，定义卷积积
\begin{equation}\tag{2.1}\label{L19eq21}(f_1*f_2)(g):=\int_{G(F)}f_1(x)f_2(x^{-1}g)\,d\mu(x).\end{equation}
于是$f_1*f_2$仍为Schwartz函数。
\sourcepage{160}
令$\mathcal H:=\mathcal H(G(F),G(\OO_F))$为$G(F)$上所有双$G(\OO_F)$不变的Schwartz函数组成的空间。卷积积$*$赋予$\mathcal H$一个代数结构，其单位元为$G(\OO_F)$的特征函数$1_{G(\OO_F)}$。$\mathcal H$称为\textbf{球面Hecke代数}\znote{球面 Hecke 代数 $\mathcal H(G(F),G(\OO_F))$ 由双 $G(\OO_F)$ 不变的紧支撑局部常值函数组成，乘法是卷积，单位元是 $G(\OO_F)$ 的特征函数。它是 $G(F)$ 的球面表示（含 $G(\OO_F)$-不变向量的不可约表示）所对应的交换 Hecke 代数：未分歧表示正对应到 $\mathcal H$ 到 $\C$ 的代数同态。$p$ 进群上它可视为有限 Hecke 代数的整体版本，也是 Satake 同构的出发空间。[来源：I. Satake, \textit{Theory of spherical functions on reductive algebraic groups over $p$-adic fields}, IHES Publ. Math. 18 (1963) 5--69（本讲参考文献 [S]）。]}，其中的函数称为球面函数。以$X$和$X^\vee$分别表示$T$的全部代数特征标和余特征标组成的格。有以下引理。
\originalnumber{1}
\begin{lemma}\label{L19S2O1}$T(F)/T(\OO_F)\cong X^\vee$。\end{lemma}
\begin{proof}
对每个$t\in T(F)$，存在唯一的余特征标$\xi(t)\in X^\vee$，使
\[\langle\lambda,\xi(t)\rangle=v(\lambda(t)),\qquad\lambda\in X,\]
其中$v:F\to\Z\cup\{+\infty\}$为$F$的归一化离散赋值。按定义，对$a\in F^\times$，有$\lambda\circ\xi(t)(a)=a^{\langle\lambda,\xi(t)\rangle}$。由于$T(F)\cong(F^\times)^r$，我们固定这样一个同构。写$t=(t_1,\ldots,t_r)$，其中$t_j=u_j\pi^{n_j}$，$u_j\in\OO_F^\times$。则$\xi(t)$为$a\mapsto(a^{n_1},\ldots,a^{n_r})$。于是$\xi(t)$平凡当且仅当所有$n_j=0$，亦即$t\in T(\OO_F)$。
% [Tutor Note: 原书 p160 引理 2.1 证明开头作「存在唯一的 $X^\vee\ni\xi(t):F^\times\to T(F),\ \pi\mapsto t$」，把余特征标与它在 $\pi$ 处的取值混在一起；$\xi(t)$ 应是 $X^\vee$ 中的余特征标（$G_m\to T$），译文按此处理。]
% [Tutor Note: 原书 p160 接着称「$\xi(t)$ 平凡当且仅当 $t\in\OO_F^\times$，等价于 $t_1=\cdots=t_r=1$」；$t\in\OO_F^\times$ 只是 $r=1$ 的情形，一般应为所有 $n_j=0$，即 $t\in T(\OO_F)=(\OO_F^\times)^r$，译文已改正。]

容易看出，$t\mapsto\xi(t)$给出满射$T(F)\to X^\vee$，其核为$T(\OO_F)$。
\end{proof}
记$\pi^{\mathbf n}=(\pi^{n_1},\ldots,\pi^{n_r})\in T(F)$，也就是$t=\pi^{\mathbf n}$时的$\xi(t)(\pi)$。选定正根后，记$D^+\subseteq\Z^r$为在上述坐标中对应于支配余权\znote{支配余权是满足 $\langle\alpha,\lambda\rangle\ge0$（对所有正根 $\alpha$）的余权 $\lambda\in X^\vee$。对分裂群，Cartan 分解的双陪集 $G(\OO_F)\lambda(\pi)G(\OO_F)$ 由支配余权唯一标记，因此 $\mathcal H$ 的基 $\{c_\lambda\}$ 与 $(X^\vee)^+$ 一一对应。几何上 $\lambda$ 记录极大环面在 $\pi$-进方向的位置；在对偶群一侧它对应支配权，而支配权与不可约有理表示一一对应，这正是 Satake 参数能读出表示不变量的原因。[来源：I. Satake, \textit{Theory of spherical functions on reductive algebraic groups over $p$-adic fields}, IHES Publ. Math. 18 (1963) 5--69（本讲参考文献 [S]）。]}的集合。有\textbf{Cartan分解}
\begin{equation}\tag{2.2}\label{L19eq22}G(F)=\bigsqcup_{\mathbf n\in D^+}G(\OO_F)\pi^{\mathbf n}G(\OO_F).\end{equation}
% [Tutor Note: 原书 p160 的 Cartan 分解写作对 $\mathbf n\in\Z^r_{\ge0}$ 求和；$\Z^r_{\ge0}$ 在 $X^\vee$ 的整基下只对应「非负坐标」的余权，双陪集实际由支配余权（一个 Weyl 室）标记，译文改用 $D^+$。]
另一种说法是：给对偶根系$\Phi(G,T)^\vee$选定基$\Delta^\vee$，便得到支配余权集合
\[(X^\vee)^+=\{\lambda\in X^\vee\mid\langle\alpha,\lambda\rangle\ge0,\ \alpha\in\Delta\}.\]
所选的正根同时确定Langlands对偶群$G^L$的一个Borel子群。\eqref{L19eq22}可写成
\[G(F)=\bigsqcup_{\lambda\in(X^\vee)^+}G(\OO_F)\lambda(\pi)G(\OO_F).\]
因此$G(\OO_F)\backslash G(F)/G(\OO_F)\cong(X^\vee)^+$，而$\mathcal H$具有基$\{c_\lambda\}_{\lambda\in(X^\vee)^+}$，其中$c_\lambda$为双陪集$G(\OO_F)\lambda(\pi)G(\OO_F)$的特征函数。\idx{Cartan分解}\idx{支配余权}
\originalnumber{2}
\begin{proposition}\label{L19S2O2}球面Hecke代数$\mathcal H$是交换的。\end{proposition}
\sourcepage{161}
\begin{proof}
只需确定结构常数
\[c_\lambda*c_\gamma=\sum_{\nu\in(X^\vee)^+}c^\nu_{\lambda\gamma}c_\nu.\]
为简化公式，记$K=G(\OO_F)$。按定义，对$a\in K\nu(\pi)K$，有
\begin{align*}
c^\nu_{\lambda\gamma}&=(c_\lambda*c_\gamma)(a)=\int_{G(F)}c_\lambda(x)c_\gamma(x^{-1}a)\,d\mu(x)\\
&=\mu\bigl(K\lambda(\pi)K\cap aK\gamma(\pi)^{-1}K\bigr)\\
&=\left|\bigl(K\lambda(\pi)K\cap aK\gamma(\pi)^{-1}K\bigr)/K\right|.
\end{align*}
可以选择$G$的Cartan对合$\sigma$。例如，对某个把Borel子群映入上三角矩阵群的嵌入$G\hookrightarrow\GL_n$，可取$\sigma(g)=(g^{\mathsf T})^{-1}$，其中$g^{\mathsf T}$表示转置。$\sigma$为$G$的自同构，满足：
\begin{enumerate}[label=\arabic*.]
\item $\sigma(g_1g_2)=\sigma(g_1)\sigma(g_2)$；
\item $\sigma(T)=T$，且$\sigma(t)=t^{-1}$，$t\in T$；
\item $\sigma(K)=K$，所以对任意可测集$V\subseteq G(F)$，$\mu(V)=\mu(\sigma(V))$。
\end{enumerate}
令$\widetilde a=\sigma(a)^{-1}\in K\nu(\pi)K$，则
\[\sigma\bigl(K\lambda(\pi)K\cap aK\gamma(\pi)^{-1}K\bigr)=K\lambda(\pi)^{-1}K\cap\widetilde a^{-1}K\gamma(\pi)K.\]
容易得到
\begin{align*}
c^\nu_{\lambda\gamma}&=\mu\bigl(K\lambda(\pi)K\cap aK\gamma(\pi)^{-1}K\bigr)\\
&=\mu\bigl(\sigma(K\lambda(\pi)K\cap aK\gamma(\pi)^{-1}K)\bigr)\\
&=\mu\bigl(K\lambda(\pi)^{-1}K\cap\widetilde a^{-1}K\gamma(\pi)K\bigr)\\
&=\mu\bigl(\widetilde aK\lambda(\pi)^{-1}K\cap K\gamma(\pi)K\bigr)=c^\nu_{\gamma\lambda}.
\end{align*}
\end{proof}
设$U$为由$T$正规化的一个Borel子群的幂幺根基。有以下\textbf{Iwasawa分解}：
\begin{equation}\tag{2.3}\label{L19eq23}G(F)=\bigsqcup_{\lambda\in X^\vee}G(\OO_F)\lambda(\pi)U(F).\end{equation}\idx{Iwasawa分解}
\sourcepage{162}
\section{Satake同构}\label{L19S3}\idx{Satake同构}
$F,\OO_F,G,\mathcal H$等仍如前。我们已经看到，$\mathcal H$由$c_\lambda$，$\lambda\in(X^\vee)^+$，张成。若$G=T$为秩$r$的$\OO_F$分裂环面，则$G(\OO_F)\cong(\OO_F^\times)^r$，且
\[c_\lambda*c_\gamma=c_{\lambda+\gamma}.\]
因而存在代数同构
\[\mathcal H(T(F),T(\OO_F))\xrightarrow{\sim}\C[x_1^{\pm1},\ldots,x_r^{\pm1}]\cong\C[X^\vee],\]
把$\sum C_\lambda c_\lambda$送到$\sum C_\lambda x^\lambda$。这里$x^\lambda=x_1^{\lambda_1}\cdots x_r^{\lambda_r}$，$\lambda=\lambda_1\gamma_1+\cdots+\lambda_r\gamma_r$，$\gamma_1,\ldots,\gamma_r$为余特征标格$X^\vee$的一组整基，$C_\lambda\in\C$。
% [Tutor Note: 原书 p162 称 $\gamma_1,\ldots,\gamma_r$ 是「all the simple coroots」；单余根一般只生成根格 $Q^\vee$，不构成 $X^\vee$ 的整基（如 $\GL_n$ 的 $X^\vee\cong\Z^n$ 需 $n$ 个元素而只有 $n-1$ 个单余根），译文改为 $X^\vee$ 的一组整基。]

然而，对一般的$G$，映射$c_\lambda\mapsto x^\lambda$并不是从$\mathcal H$到$\C[X^\vee]$的代数同态。Satake构造了另一个映射$\mathcal H\to\C[X^\vee]$，并证明了下面称为\textbf{Satake同构}\znote{Satake 同构断言球面 Hecke 代数同构于 Weyl 群不变的多项式环：$\mathcal H\cong\C[X^\vee]^W$。代数侧看，它把 $G(F)$ 的未分歧表示（即 $\mathcal H$ 的代数同态）与半单共轭类（即 $\C[X^\vee]^W$ 的极大理想）对应起来，因此未分歧表示由半单 $L$-参数决定，这就是「Satake 参数」。显式同构 $f\mapsto\widehat f$ 由 Fourier 变换在带状球函数上的取值给出，正是本讲末尾构造的内容。[来源：I. Satake, \textit{Theory of spherical functions on reductive algebraic groups over $p$-adic fields}, IHES Publ. Math. 18 (1963) 5--69（本讲参考文献 [S]）。]}的结果。
\originalnumber{1}
\begin{theorem}\label{L19S3O1}存在代数同构$\mathcal H\xrightarrow{\sim}\C[X^\vee]^W$，其中$W$为Weyl群。\end{theorem}
\originalnumber{2}
\begin{note}[L19S3O2]
实际上，可将$\C[X^\vee]$看作$X\otimes\C$上的指数多项式。在同构$\C[X^\vee]\cong\C[x_1^{\pm1},\ldots,x_r^{\pm1}]$下，$x_i$对应的函数为
\[x_i(\alpha)=e^{-\langle\alpha,\gamma_i\rangle},\qquad\alpha\in X\otimes\C,\]
其中$\gamma_1,\ldots,\gamma_r$为上述整基。换言之，同构$\mathcal H(T(F),T(\OO_F))\xrightarrow{\sim}\C[X^\vee]$由$c_\lambda\mapsto e^\lambda$给出。因此Weyl群$W$自然作用于$\C[X^\vee]$。
% [Tutor Note: 原书 p162 注记 3.2 作 $x_i(\alpha)=e^{\langle\alpha,\gamma_i\rangle}$；与本讲 (3.2)、(3.3) 采用的 $e^{-\langle s,\lambda\rangle}$ 约定相差一个负号，译文统一取 $x_i(\alpha)=e^{-\langle\alpha,\gamma_i\rangle}$。]
\end{note}
我们想简要说明\eref{定理}{L19S3O1}的证明。以下遵循\bref{L19S}{S}，先给出一个定义。
\originalnumber{3}
\begin{definition}\label{L19S3O3}
若$G(F)$上的连续复值函数$\omega$满足以下条件，则称其为$G(F)$上相对于$G(\OO_F)$的\textbf{带状球函数}\znote{带状球函数是 $G(F)$ 上双 $G(\OO_F)$ 不变、在单位元取 $1$、且是每个 $f\in\mathcal H$ 所定义卷积算子的特征函数的连续函数。它与「球面表示」一一对应：在不可约球面表示 $(\pi,V)$ 中取 $G(\OO_F)$-不变向量 $v$，矩阵元 $g\mapsto\langle\pi(g)v,v\rangle$ 就是带状球函数；反过来 $\omega$ 的 $E_f$ 给出 $\mathcal H$ 的一个特征。故 $\Omega$ 与 $\mathcal H$ 的谱（未分歧表示的 Satake 参数）是同一件事。[来源：I. Satake, \textit{Theory of spherical functions on reductive algebraic groups over $p$-adic fields}, Publ. Math. IHÉS 18 (1963) 5--69，第 1--2 节（本讲参考文献 [S]）。]}（zonal spherical function，以下简记为z.s.f.）：
\begin{enumerate}[label=(\roman*)]
\item $\omega$双$G(\OO_F)$不变；
\item $\omega(1)=1$；
\sourcepage{163}
\item 对任意$f\in\mathcal H$，$\omega$为$f$所定义积分算子的特征函数，即
\[f*\omega(g)=\int_{G(F)}f(x)\omega(x^{-1}g)\,d\mu(x)=E_f\omega(g),\qquad g\in G(F),\]
其中$E_f\in\C$。记$\widehat\omega(f)=E_f$，则
\begin{equation}\tag{3.1}\label{L19eq31}\widehat\omega(f)=\int_{G(F)}f(x)\omega(x^{-1})\,d\mu(x).\end{equation}
\end{enumerate}
\end{definition}\idx{带状球函数}
\eqref{L19eq31}中的映射$\widehat\omega$为从$\mathcal H$到$\C$的满代数同态。记$\Omega:=\Omega(G(F),G(\OO_F))$为$G(F)$上相对于$G(\OO_F)$的所有带状球函数的集合。

如何构造这样的函数？设$\chi$为$T(F)$的拟特征标，即从$T(F)$到$\C^\times$的连续同态，并满足$\chi(T(\OO_F))=1$。以$\delta(h)=|\det(\operatorname{Ad}(h)|_{\operatorname{Lie}U})|$表示相应的模特征标。借助\eqref{L19eq23}中的Iwasawa分解，$\chi$在$G(F)$上确定一个良定义的、归一化的复值函数$f_\chi$：
\[f_\chi(khu)=\delta(h)^{-1/2}\chi(h),\qquad k\in G(\OO_F),\ u\in U(F),\ h\in T(F).\]
% [Tutor Note: 原书 p163 定义 $f_\chi(khu)=\chi(h)$，未作归一化；为使其积分为带状球函数（$\omega_\chi(1)=1$ 且 $f*\omega=E_f\omega$ 成立），译文改用归一化的 $\delta(h)^{-1/2}\chi(h)$（$\delta$ 为 $T$ 的模特征标），与 Satake 的标准约定一致。]
从定义容易看出，若$g\in G(\OO_F)U(F)$，则$f_\chi(g)=1$。定义
\[\omega_\chi(g)=\int_{G(\OO_F)}\overline{f_\chi(x)}f_\chi(g^{-1}x)\,d\mu(x)=\int_{G(\OO_F)}f_\chi(g^{-1}x)\,d\mu(x).\]
则$\omega_\chi$为$G(F)$上相对于$G(\OO_F)$的带状球函数。

由\eref{引理}{L19S2O1}，这样的拟特征标$\chi$由$\chi(\pi^\lambda)$唯一确定，其中$\pi^\lambda=\lambda(\pi)$，$\lambda\in X^\vee$。因此可找到$s\in X\otimes\C$，使
\begin{equation}\tag{3.2}\label{L19eq32}\chi(\pi^\lambda)=e^{-\langle s,\lambda\rangle}.\end{equation}
这样的$s$在模$2\pi iX'$的意义下唯一，其中
\[X'=\{s\in X\otimes\C\mid\langle s,\lambda\rangle\in\Z,\ \lambda\in X^\vee\}=X.\]
% [Tutor Note: 原书 p163 作 $X\subset X'=\{s\in X\otimes\C\mid\langle s,\lambda\rangle\in\Z,\ \forall\lambda\in X^\vee\}$，只断言 $X$ 含于 $X'$；由 $X$ 与 $X^\vee$ 的配对完美，实际有 $X'=X$，译文直接写出等号。]
当$s,\chi$满足\eqref{L19eq32}时，记$\omega_s=\omega_\chi$。有以下引理。
\originalnumber{4}
\begin{lemma}\label{L19S3O4}
\begin{enumerate}[label=(\arabic*)]
\item $\omega_{-s}(g)=\omega_s(g^{-1})$。
\item 对任意$w\in W$，$\omega_{ws}=\omega_s$。
\end{enumerate}
\end{lemma}
\sourcepage{164}
第一条可以直接证明，但第二条的证明并不平凡，可见\bref{L19S}{S}。
\originalnumber{5}
\begin{theorem}\label{L19S3O5}
$G(F)$上相对于$G(\OO_F)$的所有带状球函数都可写成$\omega=\omega_s$，其中$s\in X\otimes\C$；并且$\omega_s=\omega_{s'}$当且仅当$s,s'$在$W\ltimes(2\pi iX')$的作用下等价。特别地，$\Omega$与商空间
\[(X\otimes\C)\big/\bigl(W\ltimes(2\pi iX')\bigr)\]
在复解析意义下同构。
\end{theorem}
现在，对任意$f\in\mathcal H$，定义其\textbf{Fourier变换}$\widehat f$为$X\otimes\C$上的函数
\begin{equation}\tag{3.3}\label{L19eq33}\widehat f(s):=\widehat\omega_s(f)=\int_{G(F)}f(x)\omega_s(x^{-1})\,d\mu(x).\end{equation}
在\eqref{L19eq11}所给出的Lie代数分解中，令$\mathfrak n=\bigoplus_{\alpha>0}\g_\alpha$。以$\operatorname{Ad}_{\mathfrak n}$表示$T$的伴随表示在$\mathfrak n$上的限制。定义
\[\delta:T(F)\to\R_{>0},\qquad h\longmapsto|\det(\operatorname{Ad}_{\mathfrak n}(h))|,\]
其中$|\cdot|$为赋值$v$在$F$上诱导的绝对值。$\delta$也称为$T$上的\textbf{模特征标}。定义
\[\widetilde f(h)=\delta(h)^{1/2}\int_{U(F)}f(hu)\,d\mu_U(u),\qquad h\in T(F),\]
其中$\mu_U$为$U(F)$上的Haar测度，归一化为$\mu_U(U(\OO_F))=1$。于是
\[\widehat f(s)=\sum_{\lambda\in X^\vee}\widetilde f(\pi^\lambda)e^{-\langle s,\lambda\rangle}.\]
最后，\eref{定理}{L19S3O1}中的同构显式地由$f\mapsto\widehat f$给出。\idx{模特征标}
\needspace{5\baselineskip}\section{第19讲参考文献}\label{L19S4}
\begin{references}
\refitem{L19S}{S} I. Satake, \textit{Theory of spherical functions on reductive algebraic groups over p-adic fields}, IHES Publ. Math. 18 (1963) 5--69.
\refitem{L19Sp}{S} T.A. Springer, \textit{Reductive groups}, in \textit{Automorphic forms, representations and L-functions}, Proc. Symp. Pure Math. 33, Part 1, pp. 3--27, AMS, 1979.
\end{references}
